%%%%%%%%%%%%%%%%%%%%%%%%%%%%%%%%%%%%%%%%%
%
% CMPT 435
% Lab Zero
%
%%%%%%%%%%%%%%%%%%%%%%%%%%%%%%%%%%%%%%%%%

%%%%%%%%%%%%%%%%%%%%%%%%%%%%%%%%%%%%%%%%%
% Short Sectioned Assignment
% LaTeX Template
% Version 1.0 (5/5/12)
%
% This template has been downloaded from: http://www.LaTeXTemplates.com
% Original author: % Frits Wenneker (http://www.howtotex.com)
% License: CC BY-NC-SA 3.0 (http://creativecommons.org/licenses/by-nc-sa/3.0/)
% Modified by Alan G. Labouseur  - alan@labouseur.com, and Patrick Tyler - Patrick.Tyler1@marist.edu
%
%%%%%%%%%%%%%%%%%%%%%%%%%%%%%%%%%%%%%%%%%

%----------------------------------------------------------------------------------------
%	PACKAGES AND OTHER DOCUMENT CONFIGURATIONS
%----------------------------------------------------------------------------------------

\documentclass[letterpaper, 10pt]{article} 

\usepackage[english]{babel} % English language/hyphenation
\usepackage{graphicx}
\usepackage[lined,linesnumbered,commentsnumbered]{algorithm2e}
\usepackage{listings}
\usepackage{fancyhdr} % Custom headers and footers
\pagestyle{fancyplain} % Makes all pages in the document conform to the custom headers and footers
\usepackage{lastpage}
\usepackage{url}
\usepackage{xcolor}
\usepackage{titlesec}
\usepackage{xurl}
\usepackage{multicol}
\usepackage{csvsimple}

% Stolen from https://www.overleaf.com/learn/latex/Code_listing 
\definecolor{codegreen}{rgb}{0,0.6,0}
\definecolor{codegray}{rgb}{0.5,0.5,0.5}
\definecolor{codepurple}{rgb}{0.58,0,0.82}
\definecolor{backcolour}{rgb}{0.95,0.95,0.92}

\lstdefinestyle{mystyle}{
    backgroundcolor=\color{backcolour},   
    commentstyle=\color{codegreen},
    keywordstyle=\color{magenta},
    numberstyle=\tiny\color{codegray},
    stringstyle=\color{codepurple},
    basicstyle=\ttfamily\footnotesize,
    breakatwhitespace=false,         
    breaklines=true,                 
    captionpos=b,                    
    keepspaces=true,                 
    numbers=left,                    
    numbersep=5pt,                  
    showspaces=false,                
    showstringspaces=false,
    showtabs=false,                  
    tabsize=2
}
\lstset{style=mystyle, language=c++}


\fancyhead{} % No page header - if you want one, create it in the same way as the footers below
\fancyfoot[L]{} % Empty left footer
\fancyfoot[R]{}

\renewcommand{\headrulewidth}{0pt} % Remove header underlines
\renewcommand{\footrulewidth}{0pt} % Remove footer underlines
\setlength{\headheight}{13.6pt} % Customize the height of the header

%----------------------------------------------------------------------------------------
%	TITLE SECTION
%----------------------------------------------------------------------------------------

\newcommand{\horrule}[1]{\rule{\linewidth}{#1}} % Create horizontal rule command with 1 argument of height

\title{	
   \normalfont \normalsize 
   \textsc{CMPT 308 - Spring 2024 - Dr. Labouseur} \\[10pt] % Header stuff.
   \horrule{0.5pt} \\[0.25cm] 	% Top horizontal rule
   \huge Lab 2 -- ~Cap Database\     	    % Assignment title
   \horrule{0.5pt} \\[0.25cm] 	% Bottom horizontal rule
}

\author{Patrick Tyler \\ \normalsize Patrick.Tyler1@marist.edu}

\date{\normalsize\today} 	% Today's date.

\begin{document}

\maketitle % Print the title

%----------------------------------------------------------------------------------------
%   CONTENT SECTION
%----------------------------------------------------------------------------------------

% - -- -  - -- -  - -- -  -
\section{Select * Queries}
\lstinputlisting{queries/select-people.sql}
\csvautotabular{queries/select-people.csv}
\lstinputlisting{queries/select-customers.sql}
\csvautotabular{queries/select-customers.csv}
\pagebreak
\lstinputlisting{queries/select-agents.sql}
\csvautotabular{queries/select-agents.csv}
\lstinputlisting{queries/select-product.sql}
\csvautotabular{queries/select-product.csv}
\lstinputlisting{queries/select-orders.sql}
\csvautotabular{queries/select-orders.csv}
\section{Primary key, Candidate key, and Superkey}
A super key is a set of fields (at least one) of a table which uniquely identifies each row.
A candidate key is a the minimal length of super keys which uniquely identfies each rows.
Therefore, every candidate key is a super key but not ever super key is a candidate key.
Any key which is of multiple fields is known as a composite key.
A primary key is the candidate key which is chosen to unique identify the row.
It is also very important to consider the business rules associated with the database when selecting keys.
A primary key should functionally be unique within the associated domain.
\section{Data Types in Tables}
Without data types whatever memory stored is uselesss series of 
   binary digits. 
Giving your datatypes the most strict type that is possible
   is crucial to getting information our of your data.
Imagine you have to store students as a table in your database
   let's berak down the data types that would need to be stored.
Columns and attributes such as first name and last name are fairly
   obviously strings of some length.
However, it is objectively bettter to separate these in most use cases
   into the two fields instead of something like "name" as being
   specific extracts more information.
Another field you would likely want to include grade year.
If these are students in a highschool, for example, it
   may be tempting to store their grade as a string.
This ends up being a terrible idea: when data gets entered
   it has full access to all string representation.
Should people enter a freshmen as freshmen, Freshmen, freshman,
   Freshman, freshy, first-year, etc.
It would be much better to store it has a enumeration of 
   possiblites.
In college, this representation would be less useful as 
   grade-level does not carry as much information as it does
   for high school.
It may be better to not even have grade-level stored and instead
   have it be derived from credits earned or courses takens.
\section{Relational Rules}
Making databases ensure relational rules upholds the referential integrity of the data helping to render 
   information.
It is good practice to follow these rules as well as infusing business rules into the database
   itself when possible.
\subsection{First Normal Form}
The first normal form maintains that there are no multi-valued structure within a field or fields
   relating to a semantically structured need of our business.
If a field has some internal structure that can be practically separated it should be. 
\subsection{Access Rows by Content Only}
As the databases internal structure is hidden from the user or programmers,
   accessing data using concepts such as index does not make sense.
This makes sense because in postgre, for example, all tables and columns are themselves
   stored in a database.
This means there is no concept of index as the underlying structure likely used is a set.

\subsection{Unique Rows}
Every row within a table should have a unique primary key rules of which have been previously mentioned.
This is need to differinate and assert no duplication of the same thing within the database.

\end{document}
